\documentclass[11pt,a4paper]{article}
\usepackage[utf8]{inputenc}
\usepackage[T1]{fontenc}
\usepackage{lmodern}
\usepackage{textcomp}
\usepackage{amsmath,amssymb}
\usepackage{graphicx}
\usepackage{booktabs}
\usepackage{longtable}
\usepackage{array}
\usepackage{listings}
\usepackage{xcolor}
\usepackage[margin=1in]{geometry}
\usepackage{float}
\usepackage[hidelinks,breaklinks]{hyperref}
\usepackage{microtype}
\setlength{\parskip}{4pt}
\lstset{basicstyle=\ttfamily\scriptsize,breaklines=true,breakatwhitespace=false,
  postbreak=\mbox{\textcolor{gray}{$\hookrightarrow$}\space},frame=single,
  framerule=0.3pt,rulecolor=\color{gray!60},columns=fullflexible,keepspaces=true,
  showstringspaces=false,numbers=none,xleftmargin=4pt,xrightmargin=4pt,
  literate={\_}{{\_}}1}
\graphicspath{{figures/}}
\providecommand{\tightlist}{}
\title{The Hubble Tension: A Conceptual Summary Without New Computation or Literature Search}
\author{HEP-Genesis-Agent session record \\ \small Project: ramp-features}
\date{2026-10-15}
\begin{document}
\maketitle
\noindent\textit{\small This document is an AI-generated account of an interactive session with the HEP-Genesis agent. Its text was written by a language model strictly from the session record; figures and numbers come only from tool results obtained during the session. The supplementary material and the transcript, generated alongside, document exactly what ran. Verify independently before use.}
\vspace{1em}

\begin{abstract}
This report records a short session in which the user asked for a two-paragraph explanation of the Hubble tension, followed by a written report. The assistant gave a conceptual summary of the disagreement between early-universe and late-universe determinations of the Hubble constant $H_0$. On the early-universe side it quoted the Planck cosmic microwave background (CMB) value $H_0 = 67.4 \pm 0.5$ km\,s$^{-1}$\,Mpc$^{-1}$ (arXiv:1807.06209). On the late-universe side it quoted the SH0ES Cepheid--supernova distance-ladder value $H_0 = 73.04 \pm 1.04$ km\,s$^{-1}$\,Mpc$^{-1}$ (arXiv:2112.04510). It described the discrepancy between them as about $5\sigma$. It also mentioned an intermediate tip-of-the-red-giant-branch (TRGB) calibration near $69.8$ km\,s$^{-1}$\,Mpc$^{-1}$ (arXiv:1907.05922) and a review of proposed explanations (arXiv:2103.01183). No compute, data-analysis, or literature-search tools were run. All values and citations come from the assistant's own knowledge rather than from any search or calculation in the session. The report therefore restates published values and adds no independent result.
\end{abstract}

\section{Introduction}
The user asked one question: ``What is the Hubble tension?'' The answer was to be two paragraphs, followed by a written report. The scope was explanatory. The user did not request any calculation, data fit, or systematic literature review, and none was performed. This document summarizes the explanation the assistant gave and states plainly where its content came from.

\section{Methods}
The assistant answered in quick mode. The only tool called was a report-generation tool. It received a short focus statement listing the topics to cover:
\begin{itemize}
\item the Planck CMB value;
\item the SH0ES distance-ladder value;
\item the approximately $5\sigma$ discrepancy between them;
\item the intermediate TRGB value;
\item the two classes of explanation, systematic errors and new physics.
\end{itemize}
The focus statement itself noted that no compute or literature tools had been run and that the citations came from the assistant's own knowledge.

No cosmological fits, likelihood analyses, or numerical checks were carried out. The quoted uncertainties were not combined into a significance by any calculation in the session. No figures were produced. The accompanying supplement contains the tool-call log.

\section{Results}
\subsection{Definition of the tension}
The assistant described the Hubble tension as a disagreement between two classes of measurement of the present-day expansion rate $H_0$.

The first class is the early-universe method. The standard $\Lambda$CDM model is fitted to the CMB and projected forward to today. By the assistant's account, Planck obtains $H_0 = 67.4 \pm 0.5$ km\,s$^{-1}$\,Mpc$^{-1}$ (Planck Collaboration 2018, arXiv:1807.06209). The assistant added that baryon acoustic oscillations (BAO) and Big Bang nucleosynthesis give similar values.

The second class is the late-universe distance ladder. Parallaxes and geometric anchors calibrate Cepheid variables. The Cepheids then calibrate Type Ia supernovae, which extend distance measurements into the Hubble flow. By the assistant's account, SH0ES obtains $H_0 = 73.04 \pm 1.04$ km\,s$^{-1}$\,Mpc$^{-1}$ (Riess et al. 2022, arXiv:2112.04510).

The assistant stated that the two values differ by about $5\sigma$. It added that a gap this large is too big to attribute to chance if the quoted uncertainties are correct. This significance is the assistant's statement and was not computed in the session.

\subsection{Candidate explanations}
The assistant described two broad classes of explanation and said that neither has prevailed.

The first is systematic error in the measurements. The candidates it named were:
\begin{itemize}
\item Cepheid crowding and photometry;
\item the metallicity dependence of Cepheid luminosities;
\item the standardization of Type Ia supernovae.
\end{itemize}
It noted that calibrating with TRGB stars instead of Cepheids gives intermediate values near $69.8$ km\,s$^{-1}$\,Mpc$^{-1}$ (Freedman et al. 2019, arXiv:1907.05922). In its account, this keeps the systematics debate open.

The second is physics beyond $\Lambda$CDM. The proposals it listed were early dark energy, extra relativistic species, modified recombination, and dark-sector interactions. Citing the review by Di Valentino et al. 2021 (arXiv:2103.01183), the assistant stated that most such proposals reduce the $H_0$ discrepancy only by worsening fits to other data, such as galaxy clustering ($S_8$) or BAO.

Finally, it identified several independent methods that might resolve the question:
\begin{itemize}
\item gravitational-wave standard sirens;
\item time delays in strongly lensed quasars;
\item new JWST observations of Cepheids and red giants.
\end{itemize}

\section{Discussion}
The explanation is consistent in its own terms. The early- and late-universe values are separated by roughly $5.6$ km\,s$^{-1}$\,Mpc$^{-1}$, and the TRGB value lies between them. That intermediate position is why the assistant presented the TRGB result as relevant to the systematics question.

The main limitation is that the session produced no independent evidence. None of the quoted values, uncertainties, or the $5\sigma$ significance was retrieved from a literature search or checked by calculation. The assistant said this explicitly. The values should therefore be read as the assistant's recollection of published results, not as findings verified here.

The summary also does not include the most recent measurements. The assistant itself pointed to JWST distance-ladder papers and newer BAO data as topics not covered.

\section{Conclusions}
The session produced a conceptual explanation of the Hubble tension. It frames the tension as a roughly $5\sigma$ disagreement between the Planck CMB inference of $H_0$ within $\Lambda$CDM and the SH0ES Cepheid--supernova distance ladder. In the assistant's account, the possible resolutions divide into measurement systematics and new physics, and neither is yet favored.

No new numbers or figures were generated, and the citations were not verified by search. As a possible next step, the assistant offered to search the literature for the latest $H_0$ results, including JWST distance-ladder papers and new BAO data, and to produce a revised report incorporating figures from those papers.

\section*{References}
\begin{itemize}
\item arXiv:1807.06209 --- Planck Collaboration 2018.
\item arXiv:2112.04510 --- Riess et al. 2022.
\item arXiv:1907.05922 --- Freedman et al. 2019.
\item arXiv:2103.01183 --- Di Valentino et al. 2021 (review).
\end{itemize}

\section*{Data and code availability}
All outputs of this session are collected in the project folder \href{file:///Users/nramachandra/hep-genesis-testing/projects/ramp-features}{\texttt{/\allowbreak{}Users/\allowbreak{}nramachandra/\allowbreak{}hep-genesis-testing/\allowbreak{}projects/\allowbreak{}ramp-features}}.
Full documentation of how these results were obtained is in the supplementary material (\texttt{supplement.pdf}: implementation details, the tool-call log, the research board with every finding and its evidence, figures not shown here, and the code index) and the verbatim transcript (\texttt{transcript.pdf}), generated alongside this report.


\end{document}
